% Part:sets-functions-relations
% Chapter: size-of-sets
% Section: non-enumerability-alt

\documentclass[../../../include/open-logic-section]{subfiles}

\begin{document}

\olfileid{sfr}{siz}{nen-alt}

\olsection{\printtoken{S}{nonenumerable} संच}

\begin{editorial}
  या विभागात \olref[enm-alt]{sec} मधील व्याख्या वापरून
  $\Bin^\omega$ आणि $\Pow{\Nat}$ यांची अगणनीयता सिद्ध केली आहे;
  म्हणजे $\PosInt$ पासून आच्छादन मागण्याऐवजी~$\Nat$ शी एकास-एक
  आच्छादन आवश्यक मानले आहे.
\end{editorial}

\begin{explain}
नैसर्गिक संख्यांचा संच $\Nat$ हा अनंत आहे. तो सहजपणे
!!{enumerable} देखील आहे. परंतु
\emph{!!{nonenumerable}} संच, म्हणजे !!{enumerable} नसलेले संच,
अस्तित्वात आहेत ही लक्षणीय गोष्ट आहे
(पाहा \olref[sfr][siz][enm-alt]{defn:enumerable}).

हे आश्चर्यकारक वाटू शकते. कारण $A$ हा !!{nonenumerable} आहे असे
म्हणणे म्हणजे \emph{कोणतेही} !!{bijection} $f \colon \Nat \to A$
नाही असे म्हणणे; म्हणजे~$\Nat$ मधील अनंतपणे अनेक !!{element}s~$A$
मध्ये पाठवणारे कोणतेही फलन संपूर्ण~$A$ व्यापत नाही. म्हणून $A$ हा
!!{nonenumerable} असेल, तर नैसर्गिक संख्यांपेक्षा~$A$ मध्ये ``अधिक''
!!{element}s असतात.

संच !!{nonenumerable} आहे हे सिद्ध करण्यासाठी अनुरूप !!{bijection}
अस्तित्वात असू शकत नाही हे दाखवावे लागते. याचा सर्वोत्तम मार्ग म्हणजे
~$A$ मधील !!{element}s प्रगणित करण्याचा प्रत्येक प्रयत्न किमान एक
!!{element} वगळतो हे दाखवणे; यावरून कोणतेही फलन $f\colon \Nat \to A$
हे !!{surjective} नाही असे दिसते. हे प्रस्थापित करण्याची एक सर्वसाधारण
युक्ती म्हणजे कँटर यांची \emph{कर्णरेषीय पद्धत} वापरणे. $A$ मधील
!!{element}s $x_1$, $x_2$, \dots{} असे मांडले असता, रचनेमुळे त्या
यादीत असूच शकणार नाही असा~$A$ चा आणखी एक !!{element} आपण बनवतो.

हे सर्व उदाहरणाने उत्तम समजते. म्हणून आपले पहिले उदाहरण म्हणजे $0$
आणि $1$ यांच्या सर्व अनंत चिन्हमालांचा संच~$\Bin^\omega$. (`$\Bin$'
हे द्विमान दर्शवते आणि त्याचा अर्थ फक्त दोन घटकांचा संच $\{0,1\}$
असा घेता येतो.)\oliflabeldef{sfr:card-arithmetic:card-opps:sec}{\footnote{अधिक
नेमकेपणाने, $\Bin^\omega$ हा $0$ आणि $1$ यांच्या सर्व
$\omega$-अनुक्रमांचा संच, म्हणजे $\funfromto{\omega}{\{0,1\}}$, आहे
असे ठरवायला हवे. पण याचा अर्थ
\olref[sfr][card-arithmetic][card-opps]{sec} मध्येच स्पष्ट होईल.}}{}
तरी सध्या या किंचित शिथिल मांडणीमुळे गोंधळ होऊ नये.
\end{explain}

\begin{thm}
\ollabel{thm:nonenum-bin-omega}
$\Bin^\omega$~हा !!{nonenumerable} आहे.
\end{thm}

\begin{proof}
$\Bin^\omega$ च्या एखाद्या उपसंचाचे कोणतेही प्रगणन घ्या. म्हणजे
$s_{0}$, $s_{1}$, $s_{2}$, \dots{} अशी यादी आपल्याकडे आहे, जेथे
प्रत्येक $s_n$ ही $0$ आणि~$1$ यांची अनंत चिन्हमाला आहे. या यादीतील
$n$ व्या चिन्हमालेतील $m$ वा अंक $s_n(m)$ असू द्या.
%READERNOTE{OLSIZ-008: गोठवलेल्या इंग्रजी वर्णनात s_n(m) ला m व्या चिन्हमालेतील n वा अंक म्हटले आहे; लगेचचे वाक्य आणि सारणी मात्र तो n व्या ओळीतील, म्हणजे n व्या चिन्हमालेतील, m वा अंक असल्याचे निश्चित करतात. मराठी मजकुरात ओळ-स्तंभाशी सुसंगत निर्देशांकक्रम वापरला आहे.}
मग यादीकडे सरणी म्हणून पाहता येते, ज्यात $s_n(m)$ हा $n$ व्या ओळीत
आणि $m$ व्या स्तंभात ठेवलेला आहे:
\[
\begin{array}{c|c|c|c|c|c}
& 0 & 1 & 2 & 3 & \dots \\\hline
0 & \mathbf{s_{0}(0)} & s_{0}(1) & s_{0}(2) & s_0(3) & \dots \\\hline
1 & s_{1}(0)& \mathbf{s_{1}(1)} & s_1(2) & s_1(3) & \dots \\\hline
2 & s_{2}(0)& s_{2}(1) & \mathbf{s_2(2)} & s_2(3) & \dots \\\hline
3 & s_{3}(0)& s_{3}(1) & s_3(2) & \mathbf{s_3(3)} & \dots \\\hline
\vdots & \vdots & \vdots & \vdots & \vdots & \mathbf{\ddots}
\end{array}
\]
आता या यादीत नसलेली $0$ आणि $1$ यांची $d$ ही अनंत चिन्हमाला आपण
बनवू. तिची प्रत्येक नोंद सांगून, म्हणजे सर्व $n \in \Nat$ साठी
$d(n)$ सांगून, हे करू. सहज कल्पना अशी: वरील सरणीचा कर्ण खाली वाचायचा
(त्यामुळे ``कर्णरेषीय पद्धत'' हे नाव) आणि प्रत्येक $1$ चा $0$, तर
प्रत्येक $0$ चा~$1$ करायचा.
%READERNOTE{OLSIZ-009: गोठवलेल्या इंग्रजी सूचनेत “प्रत्येक 1 चा 0” हेच परिवर्तन दोनदा आले आहे. लगेचचे प्रकरणनिहाय सूत्र परस्परपूरक दोन्ही बदल निश्चित करते: 1 चा 0 आणि 0 चा 1. मराठी मजकुरात हे दोन्ही बदल स्पष्ट केले आहेत.}
अधिक अमूर्तपणे, कर्णावरील $n$ वा !!{element}, $s_n(n)$, हा $1$ किंवा
$0$ आहे त्यानुसार $d(n)$ हा $0$ किंवा $1$ असावा अशी व्याख्या करतो;
म्हणजे:
\[
d(n) =
\begin{cases}
1 & \text{जर $s_{n}(n) = 0$ असेल}\\
0 & \text{जर $s_{n}(n) = 1$ असेल}
\end{cases}
\]
स्पष्टच, $d \in \Bin^\omega$, कारण ती $0$ आणि $1$ यांची अनंत
चिन्हमाला आहे. पण आपण $d$ अशी बनवली की कोणत्याही $n \in \Nat$
साठी $d(n) \neq s_n(n)$. म्हणजे $d$ ही $s_n$ पेक्षा तिच्या $n$ व्या
नोंदीत वेगळी आहे. म्हणून कोणत्याही $n\in \Nat$ साठी $d \neq s_n$.
त्यामुळे $d$ ही $s_0$, $s_1$, $s_2$, \dots{} या यादीत असू शकत नाही.

$\Bin^\omega$ च्या एखाद्या उपसंचाचे स्वैर प्रगणन दिले असता ते
$\Bin^\omega$ चा कोणता तरी !!{element} वगळेल हे आपण दाखवले आहे.
म्हणून $\Bin^\omega$ या संचाचे प्रगणन नाही; म्हणजे $\Bin^\omega$ हा
!!{nonenumerable} आहे.
\end{proof}

\begin{explain}
या सिद्धतापद्धतीला ``कर्णीकरण'' म्हणतात, कारण ती~$d$ ची व्याख्या
करण्यासाठी सरणीचा कर्ण वापरते. मात्र कर्णीकरणात सरणी असणे आवश्यक नाही.
प्रत्यक्षात, सरणी किंवा प्रत्यक्ष कर्ण नसतानाही हिच्याशी मिळतीजुळती कल्पना
वापरून एखादा संच !!{nonenumerable} आहे हे दाखवता येते. पुढील फलित हे
कसे करता येते ते दाखवते.
\end{explain}

\begin{thm}
\ollabel{thm:nonenum-pownat}
$\Pow{\Nat}$ हा !!{enumerable} नाही.
\end{thm}

\begin{proof}
~$\Nat$ च्या उपसंचांची प्रत्येक यादी $\Nat$ चा कोणता तरी उपसंच वगळते
हे दाखवून आपण त्याच पद्धतीने पुढे जाऊ. म्हणून $\Nat$ च्या उपसंचांची
$N_0, N_1, N_2, \ldots$ अशी एखादी यादी आहे असे समजा. $D$ या संचाची
व्याख्या पुढीलप्रमाणे करतो: $n \in D$ तेव्हा आणि केवळ तेव्हाच, जेव्हा
$n \notin N_{n}$:
\[
D = \Setabs{n \in \Nat}{n \notin N_n}
\]
स्पष्टच $D\subseteq \Nat$. पण $D$ हा यादीत असू शकत नाही. कारण
रचनेनुसार $n \in D$ तेव्हा आणि केवळ तेव्हाच, जेव्हा $n\notin N_n$;
म्हणून कोणत्याही $n \in \Nat$ साठी $D \neq N_n$.
\end{proof}

\begin{explain}
आधीच्या सिद्धतेत कर्णाचा उल्लेख नव्हता. तरी पुढीलप्रमाणे चित्र डोळ्यांसमोर
आणल्यास त्यात कर्ण आहे असे मानता येते: $N_0$, $N_1$, \dots{} हे संच
एका सरणीत लिहिल्याची कल्पना करा; $m \in N_n$ असेल तेव्हा आणि केवळ
तेव्हाच $m$ व्या स्तंभात $m$ लिहून $N_n$ हा $n$ व्या ओळीत लिहायचा.
उदाहरणार्थ, त्या यादीतील पहिले चार संच $\{0,1,2,\dots\}$,
$\{1, 3, 5, \dots\}$, $\{0,1,4\}$ आणि $\{2,3,4,\dots\}$ आहेत असे
समजा; मग आपल्या सरणीची सुरुवात अशी असेल:
\[
\begin{array}{r@{}rrrrrrr}
  N_0 = \{ & \mathbf{0}, & 1, & 2, & & & & \dots\}\\
  N_1 = \{ &  & \mathbf{1}, &  & 3, &  & 5, & \dots\}\\
  N_2 = \{ & 0, & 1, &  &  & 4\phantom{,} &  & \}\\
  N_3 = \{ &  &  & 2, & \mathbf{3}, & 4, & & \dots\}\\
  &\vdots & & & & & & \ddots\phantom{\}}
  \end{array}
\]
मग कर्णाच्या खाली जात, $n$ कर्णावर \emph{नसेल} तेव्हा आणि केवळ
तेव्हाच $n \in D$ असे ठेवून $D$ हा संच मिळतो. म्हणून वरील उदाहरणात
आपण $0$ आणि $1$ वगळू,~$2$ समाविष्ट करू,~$3$ वगळू, इत्यादी.
\end{explain}

\begin{prob}
सर्व फलने $f \colon \Nat \to \Nat$ यांचा संच स्पष्ट कर्णरेषीय
युक्तिवादाने !!{nonenumerable} आहे हे दाखवा. म्हणजे $f_1$, $f_2$,
\dots{} ही फलनांची यादी असून प्रत्येक $f_i\colon \Nat \to \Nat$
असेल, तर या यादीत नसलेले कोणते तरी $g \colon \Nat \to \Nat$ असते,
हे दाखवा.
\end{prob}

\end{document}
